\section{Literature Review}
\label{sec:literature-review}

\subsection{The Scottish Child Payment and the wider case for income transfers}
\label{sec:lit-scp}

Due to the relatively recent nature of SCP, existing literature is limited. \citeauthor{scottishgovernment2022}'s (\citeyear{scottishgovernment2022}) interim evaluation is qualitative and has no counterfactual, so it cannot separate SCP's effect from the many other things changing in Scotland and England over the same period. \citet{andersen2025} provide the closest quantitative comparison to date: a covariate-adjusted, single-aggregate OLS regression of material deprivation and food insecurity in Scotland against England, whose six-variable specification is replicated exactly in this paper as a starting point. \citet{nesom2025} use the same underlying data in a separate difference-in-differences framework, but to examine labour-supply responses to SCP, not material deprivation. Between them, these two studies leave open exactly the gap this paper is built to fill (Section~\ref{sec:introduction}).

This thin direct evidence base makes the broader cash-transfer literature useful context for whether SCP could plausibly move material hardship, not just income on paper. Unconditional and near-cash transfers elsewhere have been shown to reduce material hardship and shift spending toward child-relevant goods \citep{haushofershapiro2016,forget2011,greggwaldfogelwashbrook2006}, supporting the plausibility of a similar effect from SCP, but none is set in the UK's current welfare and cost-of-living context, and none tests a weekly, means-tested, child-specific transfer of SCP's particular design.

\subsection{Measuring child material deprivation}
\label{sec:lit-measurement}

Whatever effect SCP has on income, whether it registers as material deprivation depends on how that outcome is measured. \citet{townsend1979} established the case for measuring poverty through material conditions rather than income alone, showing that households poor on material grounds do not map neatly onto those poor on income grounds; \citet{sen1985}'s capabilities framework offers a complementary rationale, treating income as only one imperfect determinant of deprivation. This is why material deprivation, not income, is the outcome studied here. \citet{nolanwhelan1996} and \citet{whelanlaytemaitre2004} confirm the divergence holds empirically over time, while \citet{piachaud1981} shows Townsend's own threshold was itself an artefact of item selection and scoring rather than a natural discontinuity. Aggregating the ten deprivation items into a single flag is therefore a substantive modelling choice, not a neutral one. Two further strands sharpen this for the data used here: \citet{mainbradshaw2016} show UK composite indices are sensitive to exactly this kind of item and threshold choice, and \citet{nicoriciuelliot2025}, applying latent class analysis to the same HBAI item battery, find a five-class typology that the official $\geq$25 threshold cuts across rather than between. This is evidence the binary flag compresses heterogeneity that this paper's item-level decomposition is designed to recover. Measurement of the items is also imperfect: classical measurement error in a binary outcome attenuates estimated effects toward zero \citep{boundbrownmathiowetz2001}, and such error is typically non-classical and can differ systematically between groups \citep{boundkrueger1991,meyermoksullivan2015}. A concrete instance is DWP's FYE2024 questionnaire redesign (Section~\ref{sec:data}), whose potential for differential measurement across countries this paper tests directly rather than assumes away. Alongside the official DWP MDCH flag, item-level analysis lets an aggregate result be traced to, or ruled out at, individual hardships, and a \citet{benjaminihochberg1995} false discovery rate correction protects against the resulting multiple-testing risk. Two alternative composite outcome definitions are also estimated as a robustness check (Appendix~\ref{sec:appendix-composites}). This correction bounds the expected share of false positives among the tests flagged significant, a standard suited to an exploratory item-level decomposition where the priority is flagging plausible effects for further scrutiny.

\subsection{Causal identification: DiD and treatment timing}
\label{sec:lit-identification}

The identification strategy follows the DiD tradition established by \citet{cardkrueger1994}, under the parallel-trends assumption \citep{angristpischke2009}: in the absence of SCP, material deprivation in Scotland and England would have moved together. \citet{heckmanichimuratodd1998} and \citet{abadie2005} show conditioning on observables strengthens this assumption when treated and control groups differ on characteristics predicting outcome dynamics -- as Scotland and England do on housing tenure, benefit generosity and household structure -- which motivates the covariate-adjusted specifications used here, built on the same six controls as the CASE evaluation this paper partly replicates \citep{andersen2025}. SCP's own staggered rollout (under-6s from February 2021; all under-16s from November 2022) raises a second, separate identification concern documented by \citet{goodmanbacon2021} and \citet{callawaysantanna2021}: two-way fixed effects estimators can be badly biased under staggered treatment timing when treatment effects vary across cohorts or over time. The headline design here pools both phases into a single post-November-2022 indicator, following \citet{nesom2025} and \citet{andersen2025}, given limited statistical power in the small under-6 Scottish subgroup; an age-cohort staggered-timing extension is left as a background robustness check for future work, since the primary cross-fitted estimator used here is not designed for staggered adoption timing. A separate inference problem arises because Scotland is the only treated cluster among roughly ten regions: \citet{cameronmiller2015} and \citet{mackinnonwebb2018} show cluster-robust and wild-bootstrap inference both become unreliable with so few clusters and only one treated. A placebo test confirms this diagnosis directly for this data: most placebo English regions "fail" the same pre-trend test Scotland does. This is resolved by clustering at household level instead, which is independently the substantively correct level given the within-child correlation across the ten items \citep{moulton1990}.

\subsection{Doubly robust and machine-learning-augmented estimators}
\label{sec:lit-dml}

The design above still leans on OLS's linear functional form to adjust for the six CASE covariates; clustering fixes the standard errors, but if the true relationship between those covariates and deprivation is nonlinear, or involves interactions the linear model does not capture, the point estimate itself, not just its precision, may be biased. A separate estimation literature offers two alternative ways to adjust for covariates without assuming a particular functional form, each with a different, single point of failure. \citet{abadie2005}'s inverse probability weighting (IPW) estimator reweights control-group observations by the inverse of their estimated probability of being treated, so that the reweighted controls resemble the treated group on observables; it is consistent only if this propensity model is correctly specified. An outcome-regression estimator instead models the outcome directly as a function of covariates within each treatment/period cell and extrapolates the counterfactual from that model; it is consistent only if the outcome model, not the propensity model, is correctly specified. Both the propensity model and the outcome model are examples of what this literature calls a nuisance model: a component that must be estimated to recover the treatment effect but is not itself of interest. \citet{santannazhao2020} combine both nuisance models into a single doubly robust DiD estimator for repeated cross-sections that is consistent if either one is correctly specified, not necessarily both. This is directly relevant here, since the linear six-variable OLS specification being replicated could itself be misspecified in either role.

Naively substituting a flexible machine-learning model for either nuisance model reintroduces a different problem: regularised methods such as lasso or random forest are deliberately biased, by design, to control variance, and that bias can leak into the treatment-effect estimate built on top of them. \citet{chernozhukov2018} show this contamination can be removed by combining a Neyman-orthogonal moment condition, one whose sensitivity to small errors in the nuisance models is zero to first order, with cross-fitting, where the nuisance models are estimated on one part of the sample and the treatment effect is evaluated on a held-out part, so that a model's own overfitting to its training data cannot bias the estimate evaluated on data it never saw. This is what allows a flexible model like random forest to replace OLS without reintroducing bias into the causal estimate.

Two doubly robust, ML-augmented DiD constructions were tested for this paper: \citet{chang2020}'s repeated-cross-section estimator, ruled out after returning an implausibly large, tightly estimated effect consistent with propensity-model collapse, and \citet{zimmert2020}'s efficient-score, doubly robust estimator, adopted instead and implemented via \texttt{causalweight::didDML()} under both lasso and random forest nuisance models (Section~\ref{sec:dml}).